%%
%% Commands for TeXCount
%TC:macro \cite [option:text,text]
%TC:macro \citep [option:text,text]
%TC:macro \citet [option:text,text]
%TC:envir table 0 1
%TC:envir table* 0 1
%TC:envir tabular [ignore] word
%TC:envir displaymath 0 word
%TC:envir math 0 word
%TC:envir comment 0 0
%%
\documentclass[sigconf, nonacm]{acmart}

\usepackage{booktabs}
\usepackage{multirow}
\usepackage{rotating}
\usepackage{multirow}
\usepackage{graphicx}
\usepackage{float}
\usepackage{placeins}



\renewcommand{\arraystretch}{1.15}

\AtBeginDocument{%
  \providecommand\BibTeX{{%
    Bib\TeX}}}

%% Remove ACM-specific elements for course use
\settopmatter{printacmref=false}
\renewcommand\footnotetextcopyrightpermission[1]{}
\pagestyle{plain}

\begin{document}

\title{Cross-Dataset Generalizability of Sarcasm Detection: Encoder Fine-Tuning vs.\ LLM Prompting}

\author{Shiyi Liu}
\email{lshiyi25@vt.edu}
\affiliation{%
  \institution{Virginia Tech}
  \department{Department of Computer Science}
  \city{Alexandria}
  \country{USA}
}

\author{Christopher Williams}
\email{cdw24@vt.edu}
\affiliation{%
  \institution{Virginia Tech}
  \department{Department of Computer Science}
  \city{Alexandria}
  \country{USA}
}

\renewcommand{\shortauthors}{Liu and Williams}

\maketitle

% ============================================================
\section{Project Idea}
% ============================================================

Sarcasm detection persists as a challenge in Natural Language Processing (NLP) due to the implied meaning, context, and speaker intent that is not explicitly communicated. Sarcasm is pervasive and a constant part of human communication, both in person and online, and failing to detect this with NLP can miss the inversion of a statement, harming downstream tasks and operations like sentiment analysis and opinion mining.

Our project will investigate how well Large Language Models (LLMs) handle sarcasm detection compared to fine-tuned encoder-based transformers. Specifically, we aim to replicate the cross-dataset evaluation framework from \cite{jang2024generalizablesarcasmdetectionjust}, which fine-tunes models such as BERT and RoBERTa on multiple sarcasm datasets, and compare those results against zero-shot and few-shot prompting with LLMs such as GPT-5 and Claude Opus 4.6. This comparison directly tests whether the scale and broad pretraining of generative models can overcome the cross-dataset generalization gap that the original paper identified as a core limitation of fine-tuned approaches.

% ============================================================
\section{Motivation}
% ============================================================

Sarcasm detection is a necessary upstream task for sentiment analysis and opinion mining, and current approaches leave a clear gap between two lines of research that our project aims to bridge.

Recent work by Jang and Frassinelli~\cite{jang2024generalizablesarcasmdetectionjust} demonstrated that fine-tuned encoder-based models struggle to generalize across sarcasm datasets with different styles, domains, and modalities, with cross-dataset F-scores dropping heavily from intra-dataset scores. This behavior shows that these models (BERT, RoBERTa, DeBERTa) learn dataset-specific cues rather than a general understanding of sarcasm and its cues.

Separately, Zhang et al.~\cite{zhang2024sarcasmbench} benchmarked eleven LLMs on sarcasm detection and found that even GPT-4 (the strongest performer in their study) underperformed supervised fine-tuned baselines in most settings. However, their evaluation did not use the cross-dataset framework from Jang and Frassinelli, leaving open the question of whether LLMs might exhibit better \textit{generalization} across sarcasm styles, domains, and modalities.

Using both fine-tuned encoders and prompted LLMs under the same cross-dataset setup, we can directly compare generalizability, seeing if the breadth of LLM pretraining compensates for the lack of task-specific fine-tuning when sarcasm shifts in appearance.

% ============================================================
\section{Datasets}
% ============================================================

We plan to use the four core datasets from Jang and Frassinelli~\cite{jang2024generalizablesarcasmdetectionjust} to enable direct comparison with their results, and include two additional datasets to broaden the range of sarcasm styles and domains evaluated. Table~\ref{tab:datasets} summarizes the datasets.

\begin{table*}[t]
\centering
\caption{Overview of datasets used in this project.}
\label{tab:datasets}
\begin{tabular}{llrlll}
\toprule
\textbf{Dataset} & \textbf{Domain} & \textbf{Size} &
\textbf{Labels} & \textbf{Context} & \textbf{Source} \\
\midrule
CSC & Simulated conversation & 7,040 & Author + Third-party &
Yes & \cite{jang2024generalizablesarcasmdetectionjust} \\
MUStARD & TV series dialogue & 690 & Third-party &
Yes & \cite{mustard} \\
SC V2 & Online debate forums & 9,386 & Third-party &
No & \cite{sarcasmcorpusv2} \\
iSarcasmEval & Twitter posts & 3,801 & Author &
No & \cite{isarcasmeval} \\
\midrule
SARC & Reddit comments & 1.3M sarcastic & Author (/s tag) &
Yes & \cite{khodak2018largeselfannotatedcorpussarcasm} \\
SAD & Twitter posts & 2,340 & Manual annotation &
No & \cite{painter2022sad} \\
\bottomrule
\end{tabular}
\end{table*}

We add SARC because it is the largest publicly available sarcasm corpus and uses self-annotation via the \texttt{/s} tag, a labeling method discussed but not tested in the original paper~\cite{khodak2018largeselfannotatedcorpussarcasm}. SAD provides an additional Twitter-based dataset with manual annotations, offering a point of comparison against iSarcasmEval's author-labeled Twitter data. Together, the six datasets span conversational, social media, debate forum, and scripted dialogue domains, with both author and third-party labeling.

% ============================================================
\section{Examples from Dataset}
% ============================================================

Please see Table~\ref{tab:examples} for examples from our datasets.

\begin{table*}[t]
\centering
\caption{Example sarcastic instances from each dataset.}
\label{tab:examples}
\small
\begin{tabular}{p{2cm}p{13.5cm}}
\toprule
\textbf{Dataset} & \textbf{Example} \\
\midrule
CSC &
\textbf{Context:} Steve gives you a watering can on your birthday
while smiling at you with a strange expression. But you don't even
have a single plant. \newline
\textbf{Response:} Maybe I will use it as an outside shower. \\
\midrule
MUStARD &
\textbf{Context:} ``How do I look?'', ``Could you be more
specific?'', ``Can you tell I'm perspiring a little?'' \newline
\textbf{Response:} No. The dark crescent-shaped patterns under
your arms conceal it nicely. \\
\midrule
SC V2 &
Aaahhh, so just not accomplishing our goals and going home is not
defeat, there has to be paperwork for it to be a defeat. Gotcha. \\
\midrule
iSarcasmEval &
Imagine going to university for 4 years when you could just follow
Elon Musk on Twitter for free. \\
\midrule
SARC &
\textbf{Parent:} Finally, the bankers have a voice in Washington!
\newline
\textbf{Comment:} Yeah, it's been really tough for them. /s \\
\midrule
SAD &
Oh great, another Monday. Just what I needed. \#sarcasm \\
\bottomrule
\end{tabular}
\end{table*}

% ============================================================
\section{NLP Approaches}
% ============================================================

\subsection{Approach 1: Fine-Tuned Encoder Models}

Following \cite{jang2024generalizablesarcasmdetectionjust}'s approach, we fine-tune three encoder-based transformers (BERT, RoBERTa, DeBERTa) for binary sarcasm classification. Each model is fine-tuned on each dataset individually using the HuggingFace Trainer API, with true 5-fold cross-validation (80/20 split, not 64/16--20 for test) and multiple random seeds. We will evaluate both intra-dataset performance (train and test on the same dataset) and cross-dataset performance (train on one, test on others). This replication serves as a baseline and allows us to verify the original paper's findings before comparing against LLM-based approaches.

\subsection{Approach 2: Zero-Shot and Few-Shot LLM Prompting}

We will evaluate LLMs (e.g., GPT-5, Claude Opus 4.6, Gemini 3.0, Grok 4), subject to availability and API access, on the same datasets using zero-shot and few-shot prompting. In the zero-shot setting, the model receives only a task description and the target text. In the few-shot setting, we provide 2--4 labeled examples (balanced between sarcastic and non-sarcastic) sampled from the training split of one dataset before asking the model to classify instances from the same dataset, then another round classifying instances from a different dataset. This mimics the cross-dataset setup from Approach~1, enabling direct comparison. By testing if LLMs generalize better across sarcasm styles and domains without any task-specific fine-tuning, we address the gap between the findings of~\cite{jang2024generalizablesarcasmdetectionjust} and~\cite{zhang2024sarcasmbench}.

% ============================================================
\section{Metrics}
% ============================================================

To enable direct comparison with Jang and
Frassinelli~\cite{jang2024generalizablesarcasmdetectionjust}, we use \textbf{F1-score} as our primary evaluation metric. F1 is preferred over accuracy for this task because sarcasm datasets vary in class balance, and accuracy can be misleading when a model simply defaults to the majority class.

We additionally report \textbf{precision} and \textbf{recall} separately, as they reveal distinct failure modes: low precision indicates the model over-predicts sarcasm (false positives), while low recall indicates it misses sarcastic instances (false negatives). This distinction is relevant for cross-dataset evaluation, where a model trained on debate-forum corpora (e.g., SC V2) may either over-trigger or under-trigger on subtler conversational sarcasm (e.g., CSC).

For each approach, we report these metrics in two settings: \textbf{intra-dataset} (train and test on the same dataset) and \textbf{cross-dataset} (train on one dataset, test on another). The cross-dataset results will form a view similar to Table~4 in the original paper, allowing us to directly compare the generalization profiles of fine-tuned encoders and prompted LLMs. We also compute the \textbf{average cross-dataset F1} per training source as a summary measure of generalizability, following the original paper's analysis. Finally, we will conduct qualitative error analysis to better understand model behavior across datasets.


\section{Methodology}

% ============================================================
\section{Results}
% ============================================================

\subsection{Encoders}

In table~\ref{tab:encoder_results} we compare the results of all encoders on intra dataset training and testing, and then cross dataset training and testing. 

%% longtable is incompatible with sigconf's two-column layout.
%% Solution: use table* (full-width float) with plain tabular, split into two floats.
%% landscape is still used via pdflscape to give the wide table enough horizontal room.

% --- Table 3a: BERT, DeBERTa, DistilBERT ---
\begin{table*}[htbp]
\centering
\caption{F1 scores for all encoder models including CSC label type (+A = author, default = third-party) and context (+C) variants. Intra F1 = train and test on the same dataset. Avg Cross F1 = mean over all valid cross-dataset transfers (same-source-corpus pairs excluded). Values rounded to two decimal places.}
\label{tab:encoder_results_full}
\resizebox{0.9\textwidth}{!}{
\begin{tabular}{l|c|cccccccccc|c}
\toprule
Train Dataset & Intra F1 & CSC & CSC+A & CSC+A+C & CSC+C & iSarcasm & MUStARD & MUS+C & News & SARC & SarcV2 & Avg Cross F1 \\
\midrule
\multicolumn{13}{l}{\textbf{BERT}} \\
CSC & 0.39 & -- & -- & -- & -- & 0.28 & 0.24 & 0.15 & 0.41 & 0.28 & 0.23 & 0.26 \\
CSC+A & 0.96 & -- & -- & -- & -- & 0.32 & 0.51 & 0.39 & 0.56 & 0.38 & 0.34 & 0.42 \\
CSC+A+C & 0.99 & -- & -- & -- & -- & 0.33 & 0.49 & 0.30 & 0.55 & 0.35 & 0.36 & 0.40 \\
CSC+C & 0.33 & -- & -- & -- & -- & 0.08 & 0.16 & 0.11 & 0.15 & 0.08 & 0.14 & 0.12 \\
iSarcasm & 0.25 & 0.17 & 0.17 & 0.00 & 0.01 & -- & 0.10 & 0.06 & 0.00 & 0.22 & 0.26 & 0.11 \\
MUStARD & 0.61 & 0.27 & 0.30 & 0.31 & 0.30 & 0.40 & -- & -- & 0.65 & 0.61 & 0.65 & 0.44 \\
MUS+C & 0.67 & 0.30 & 0.32 & 0.27 & 0.28 & 0.40 & -- & -- & 0.65 & 0.66 & 0.62 & 0.44 \\
News & 0.92 & 0.13 & 0.13 & 0.00 & 0.01 & 0.22 & 0.00 & 0.00 & -- & 0.17 & 0.07 & 0.08 \\
SARC & 0.73 & 0.30 & 0.29 & 0.28 & 0.28 & 0.42 & 0.53 & 0.43 & 0.38 & -- & 0.44 & 0.37 \\
SarcV2 & 0.81 & 0.33 & 0.36 & 0.37 & 0.35 & 0.43 & 0.60 & 0.47 & 0.64 & 0.60 & -- & 0.46 \\
\midrule
\multicolumn{13}{l}{\textbf{DeBERTa}} \\
CSC & 0.38 & -- & -- & -- & -- & 0.34 & 0.41 & 0.20 & 0.40 & 0.33 & 0.39 & 0.35 \\
CSC+A & 0.95 & -- & -- & -- & -- & 0.40 & 0.64 & 0.47 & 0.55 & 0.42 & 0.52 & 0.50 \\
CSC+A+C & 0.91 & -- & -- & -- & -- & 0.39 & 0.67 & 0.30 & 0.58 & 0.59 & 0.63 & 0.53 \\
CSC+C & 0.39 & -- & -- & -- & -- & 0.41 & 0.61 & 0.30 & 0.61 & 0.57 & 0.55 & 0.51 \\
iSarcasm & 0.43 & 0.28 & 0.31 & 0.10 & 0.09 & -- & 0.47 & 0.39 & 0.17 & 0.45 & 0.56 & 0.31 \\
MUStARD & 0.66 & 0.30 & 0.31 & 0.31 & 0.30 & 0.40 & -- & -- & 0.65 & 0.67 & 0.67 & 0.45 \\
MUS+C & 0.66 & 0.30 & 0.31 & 0.31 & 0.30 & 0.40 & -- & -- & 0.65 & 0.67 & 0.67 & 0.45 \\
News & 0.93 & 0.13 & 0.15 & 0.02 & 0.01 & 0.24 & 0.16 & 0.06 & -- & 0.20 & 0.10 & 0.12 \\
SARC & 0.75 & 0.29 & 0.30 & 0.23 & 0.24 & 0.46 & 0.57 & 0.45 & 0.39 & -- & 0.40 & 0.37 \\
SarcV2 & 0.83 & 0.34 & 0.37 & 0.36 & 0.33 & 0.43 & 0.60 & 0.61 & 0.62 & 0.61 & -- & 0.47 \\
\midrule
\multicolumn{13}{l}{\textbf{DistilBERT}} \\
CSC & 0.26 & -- & -- & -- & -- & 0.11 & 0.21 & 0.19 & 0.21 & 0.19 & 0.15 & 0.18 \\
CSC+A & 0.98 & -- & -- & -- & -- & 0.32 & 0.46 & 0.37 & 0.60 & 0.42 & 0.38 & 0.42 \\
CSC+A+C & 0.97 & -- & -- & -- & -- & 0.39 & 0.61 & 0.22 & 0.62 & 0.56 & 0.49 & 0.48 \\
CSC+C & 0.38 & -- & -- & -- & -- & 0.20 & 0.29 & 0.20 & 0.37 & 0.21 & 0.21 & 0.25 \\
iSarcasm & 0.16 & 0.16 & 0.17 & 0.01 & 0.01 & -- & 0.27 & 0.11 & 0.00 & 0.19 & 0.24 & 0.13 \\
MUStARD & 0.48 & 0.28 & 0.31 & 0.31 & 0.30 & 0.40 & -- & -- & 0.66 & 0.60 & 0.64 & 0.44 \\
MUS+C & 0.70 & 0.30 & 0.33 & 0.34 & 0.30 & 0.40 & -- & -- & 0.65 & 0.66 & 0.66 & 0.45 \\
News & 0.92 & 0.10 & 0.11 & 0.04 & 0.03 & 0.20 & 0.00 & 0.00 & -- & 0.18 & 0.13 & 0.09 \\
SARC & 0.72 & 0.29 & 0.29 & 0.27 & 0.27 & 0.36 & 0.46 & 0.44 & 0.38 & -- & 0.37 & 0.35 \\
SarcV2 & 0.80 & 0.32 & 0.34 & 0.34 & 0.32 & 0.41 & 0.64 & 0.62 & 0.59 & 0.61 & -- & 0.47 \\
\midrule
\multicolumn{13}{l}{\textbf{ELECTRA}} \\
CSC & 0.29 & -- & -- & -- & -- & 0.16 & 0.16 & 0.00 & 0.24 & 0.17 & 0.09 & 0.14 \\
CSC+A & 0.97 & -- & -- & -- & -- & 0.36 & 0.60 & 0.31 & 0.59 & 0.42 & 0.52 & 0.47 \\
CSC+A+C & 0.96 & -- & -- & -- & -- & 0.31 & 0.51 & 0.31 & 0.43 & 0.42 & 0.47 & 0.41 \\
CSC+C & 0.14 & -- & -- & -- & -- & 0.02 & 0.00 & 0.00 & 0.00 & 0.00 & 0.00 & 0.01 \\
iSarcasm & 0.35 & 0.23 & 0.23 & 0.02 & 0.02 & -- & 0.43 & 0.39 & 0.01 & 0.23 & 0.43 & 0.22 \\
MUStARD & 0.56 & 0.29 & 0.30 & 0.33 & 0.32 & 0.41 & -- & -- & 0.65 & 0.61 & 0.68 & 0.45 \\
MUS+C & 0.63 & 0.30 & 0.32 & 0.32 & 0.30 & 0.40 & -- & -- & 0.65 & 0.66 & 0.64 & 0.45 \\
News & 0.94 & 0.09 & 0.07 & 0.01 & 0.01 & 0.26 & 0.06 & 0.10 & -- & 0.12 & 0.07 & 0.09 \\
SARC & 0.74 & 0.28 & 0.29 & 0.20 & 0.18 & 0.44 & 0.52 & 0.32 & 0.41 & -- & 0.46 & 0.35 \\
SarcV2 & 0.82 & 0.34 & 0.36 & 0.35 & 0.34 & 0.46 & 0.66 & 0.62 & 0.58 & 0.62 & -- & 0.48 \\
\midrule
\multicolumn{13}{l}{\textbf{RoBERTa}} \\
CSC & 0.08 & -- & -- & -- & -- & 0.04 & 0.00 & 0.00 & 0.00 & 0.06 & 0.01 & 0.02 \\
CSC+A & 0.45 & -- & -- & -- & -- & 0.36 & 0.05 & 0.29 & 0.24 & 0.36 & 0.38 & 0.28 \\
CSC+A+C & 0.00 & -- & -- & -- & -- & 0.00 & 0.00 & 0.00 & 0.00 & 0.00 & 0.00 & 0.00 \\
CSC+C & 0.00 & -- & -- & -- & -- & 0.00 & 0.00 & 0.00 & 0.00 & 0.00 & 0.00 & 0.00 \\
iSarcasm & 0.00 & 0.00 & 0.00 & 0.00 & 0.00 & -- & 0.00 & 0.00 & 0.00 & 0.00 & 0.00 & 0.00 \\
MUStARD & 0.41 & 0.20 & 0.22 & 0.32 & 0.31 & 0.37 & -- & -- & 0.62 & 0.43 & 0.64 & 0.39 \\
MUS+C & 0.00 & 0.00 & 0.00 & 0.00 & 0.00 & 0.00 & -- & -- & 0.00 & 0.00 & 0.00 & 0.00 \\
News & 0.71 & 0.09 & 0.11 & 0.24 & 0.23 & 0.21 & 0.06 & 0.16 & -- & 0.18 & 0.30 & 0.18 \\
SARC & 0.64 & 0.27 & 0.28 & 0.10 & 0.08 & 0.34 & 0.52 & 0.29 & 0.41 & -- & 0.32 & 0.29 \\
SarcV2 & 0.68 & 0.30 & 0.31 & 0.31 & 0.30 & 0.39 & 0.63 & 0.64 & 0.56 & 0.66 & -- & 0.46 \\
\bottomrule
\end{tabular}
}
\end{table*}

\subsection{LLMs}

\begin{table*}[htbp]
\centering
\caption{F1 scores for LLM prompting models. Intra F1 refers to performance when the train and test datasets match. Avg Cross F1 is the average across all non-matching datasets.}
\label{tab:llm_results}
\resizebox{\textwidth}{!}{
\begin{tabular}{l|c|ccccccc}
\toprule
Train Dataset & Intra F1 & CSC & iSarcasm & MUStARD & News & SARC & SarcV2 & Avg Cross F1 \\
\midrule

\multicolumn{9}{l}{\textbf{GPT (Zero-shot)}} \\
CSC       & 0.20 & -- & 0.45 & 0.73 & 0.74 & 0.66 & 0.78 & \textbf{0.67} \\
iSarcasm  & 0.48 & 0.10 & -- & 0.75 & 0.71 & 0.00 & 0.00 & 0.31 \\
MUStARD   & 0.73 & 0.00 & 0.48 & -- & 0.73 & 0.62 & 0.77 & 0.52 \\
News      & 0.76 & 0.11 & 0.48 & 0.75 & -- & 0.61 & 0.00 & 0.39 \\
SARC      & 0.64 & 0.00 & 0.38 & 0.75 & 0.65 & -- & 0.77 & 0.51 \\
SarcV2    & 0.73 & 0.12 & 0.53 & 0.68 & 0.68 & 0.62 & -- & 0.53 \\

\midrule
\multicolumn{9}{l}{\textbf{GPT (Few-shot)}} \\
CSC       & 0.25 & -- & 0.50 & 0.75 & 0.68 & 0.66 & 0.81 & \textbf{0.68} \\
iSarcasm  & 0.43 & 0.22 & -- & 0.60 & 0.74 & 0.00 & 0.00 & 0.31 \\
MUStARD   & 0.68 & 0.18 & 0.50 & -- & 0.62 & 0.55 & 0.80 & 0.53 \\
News      & 0.81 & 0.11 & 0.48 & 0.67 & -- & 0.58 & 0.00 & 0.37 \\
SARC      & 0.62 & 0.00 & 0.35 & 0.72 & 0.56 & -- & 0.76 & 0.48 \\
SarcV2    & 0.78 & 0.14 & 0.43 & 0.58 & 0.51 & 0.60 & -- & 0.45 \\

\midrule
\multicolumn{9}{l}{\textbf{Kimi-K2.6 (Zero-shot)}} \\
CSC       & 0.21 & -- & 0.55 & 0.82 & 0.89 & 0.72 & 0.81 & \textbf{0.76} \\
iSarcasm  & 0.40 & 0.11 & -- & 0.76 & 0.88 & 0.74 & 0.85 & 0.67 \\
MUStARD   & 0.79 & 0.10 & 0.50 & -- & 0.89 & 0.72 & 0.80 & 0.60 \\
News      & 0.88 & 0.18 & 0.45 & 0.74 & -- & 0.71 & 0.80 & 0.58 \\
SARC      & 0.74 & 0.22 & 0.52 & 0.84 & 0.92 & -- & 0.77 & 0.65 \\
SarcV2    & 0.79 & 0.22 & 0.50 & 0.79 & 0.85 & 0.70 & -- & 0.61 \\

\midrule
\multicolumn{9}{l}{\textbf{Kimi-K2.6 (Few-shot)}} \\
CSC       & 0.08 & -- & 0.45 & 0.81 & 0.86 & 0.70 & 0.87 & \textbf{0.74} \\
iSarcasm  & 0.45 & 0.25 & -- & 0.79 & 0.85 & 0.74 & 0.81 & 0.69 \\
MUStARD   & 0.74 & 0.21 & 0.58 & -- & 0.86 & 0.69 & 0.87 & 0.64 \\
News      & 0.88 & 0.11 & 0.48 & 0.72 & -- & 0.70 & 0.84 & 0.57 \\
SARC      & 0.74 & 0.09 & 0.48 & 0.81 & 0.77 & -- & 0.78 & 0.59 \\
SarcV2    & 0.84 & 0.11 & 0.59 & 0.75 & 0.85 & 0.69 & -- & 0.60 \\

\bottomrule
\end{tabular}
}
\end{table*}





\FloatBarrier 
\bibliographystyle{ACM-Reference-Format}
\bibliography{references}

\end{document}